\documentclass[11pt,a4paper]{article}

% --- Core arXiv-compatible packages ---
\usepackage[margin=1in]{geometry}
\usepackage{amsmath,amssymb,amsthm}
\usepackage{graphicx}
\usepackage{booktabs}
\usepackage{array}
\usepackage{multirow}
\usepackage{xcolor}
\usepackage{hyperref}
\usepackage{url}
\usepackage{algorithm}
\usepackage{algpseudocode}
\usepackage{enumitem}
\usepackage{caption}

\hypersetup{
    colorlinks=true,
    linkcolor=blue!60!black,
    citecolor=blue!60!black,
    urlcolor=blue!60!black,
}

% --- Theorem-like environments ---
\theoremstyle{definition}
\newtheorem{definition}{Definition}[section]
\newtheorem{assumption}{Assumption}[section]

\theoremstyle{plain}
\newtheorem{theorem}{Theorem}[section]
\newtheorem{proposition}[theorem]{Proposition}
\newtheorem{lemma}[theorem]{Lemma}
\newtheorem{remark}[theorem]{Remark}

% --- Math shortcuts ---
\newcommand{\R}{\mathbb{R}}
\newcommand{\Z}{\mathbb{Z}}
\newcommand{\E}{\mathbb{E}}
\newcommand{\norm}[1]{\left\lVert#1\right\rVert}
\newcommand{\abs}[1]{\left\lvert#1\right\rvert}
\DeclareMathOperator{\argmin}{arg\,min}
\DeclareMathOperator{\argmax}{arg\,max}
\DeclareMathOperator{\tr}{tr}
\DeclareMathOperator{\rank}{rank}
\DeclareMathOperator{\diag}{diag}
\DeclareMathOperator{\sign}{sign}

% ----------------------------------------------------------------------
\title{\textbf{KakeyaLattice: D4 Nested-Lattice KV-Cache Compression}\\
  \textbf{with Kakeya-Style Discrete Codebooks}}
\author{Allen Li \\[2pt]
  \small Individual researcher \\
  \small Email: \href{mailto:AllenL329@gmail.com}{AllenL329@gmail.com}}
\date{April 23, 2026}
% ----------------------------------------------------------------------

\begin{document}
\maketitle

\begin{abstract}
We present \textsc{KakeyaLattice}~v1.4, a KV-cache compression codec for
transformer-based large language models that is, to our knowledge, the
first Kakeya-inspired discrete-codebook construction to beat the
production TurboQuant~\cite{turboquant-vllm} codec head-to-head on
real vLLM inference across four open-source model families
(Qwen3-4B, DeepSeek-R1-Distill-Qwen-1.5B, Gemma-4-E4B, GLM-4-9B-Chat).

The codec is a Zamir--Feder nested lattice quantiser built on the $D_4$
root lattice, wrapped in the full TurboQuant engineering stack
(Sylvester--Hadamard rotation, per-vector adaptive scale, unit-norm
factorisation, matched bit budget). The $D_4$ structure contributes a
shaping gain of approximately $+0.37\,$dB over scalar $\Z^4$
quantisation; our theoretical prediction that this translates to a
K-MSE ratio of $\approx 0.92\times$ versus TurboQuant matches the
measured ratio of $0.911\times$ on live vLLM to within $1\%$. On 12
head-to-head pairs (4 models $\times$ 3 matched bit points) \textsc{KakeyaLattice}
wins universally on K-MSE (12/12, $10$--$36\%$ better), wins $9/12$ on
$\abs{\Delta\mathrm{ppl}}$ (4/4 at the aggressive $\sim\!3.6\times$
compression point, $1.8\times$--$5.3\times$ better), and reproduces the
same shaping gain across every model.  At $128$k context the codec
shrinks the K+V cache by $2.5$--$2.8\times$ at balanced quality and by
$1.66$--$1.73\times$ at near-lossless quality; the Q=152 near-lossless
point has $\abs{\Delta\mathrm{ppl}} < 1\%$ and top-1 pair agreement
$\geq 98\%$ on every tested model.

We place this construction inside the Kakeya--Brascamp--Lieb--Tropp
dependency chain --- Guth's multilinear Kakeya endpoint~\cite{guth-endpoint}
$\to$ Bennett--Carbery--Christ--Tao Brascamp--Lieb
equivalence~\cite{bennett-carbery-christ-tao} $\to$ Tropp's matrix
concentration~\cite{tropp-matrix-chernoff} $\to$ Halko--Martinsson--Tropp
randomised-algebra bounds~\cite{halko-martinsson-tropp} --- whose deep
end is used by Wang and Zahl~\cite{wang-zahl} in their recent
resolution of the three-dimensional Euclidean Kakeya conjecture. The
$D_4$ nested-lattice construction is a literature-traceable shallow
instance of the same Kakeya-style toolchain: $D_4$ realises the
densest packing lattice in $\R^4$, and the lattice's Voronoi cells
form a Kakeya-like direction-covering system for quantised residuals.
A research trail of three independent ``Dvir-to-Euclidean
bridges''~\cite{dvir-finite-field} --- Guth--Katz polynomial
partitioning, naive Zamir--Feder $D_4$, and empirical non-Gaussian
shaping --- is summarised together with the measurements that ruled
each out, isolating the combination ``$D_4$ lattice + TurboQuant
engineering stack'' as the first winning construction.

We position \textsc{KakeyaLattice}~v1.4 as an $\approx 8\%$ improvement to the
upstream TurboQuant kernel rather than a standalone codec, because
(i) the $D_4$ structure is a drop-in replacement for the
int$\cdot$per-coord quantiser inside TurboQuant's existing Hadamard +
qmax wrapper, and (ii) the engineering stack contributes $\sim$1000x
more to the measured quality than the lattice choice, dwarfing any
standalone Kakeya-inspired construction. The measured gain is
architecture-independent (same shaping factor on four different
transformer families) and carries zero quality cost at the aggressive
compression point.  All code, multi-model measurements, raw
per-passage JSON, and snapshot-mode vLLM harness are released under
Apache-2.0 at
\url{https://github.com/FluffyAIcode/LLM-KV--Cache-compress} at tag
\texttt{v1.4}.
\end{abstract}

\section{Introduction}
\label{sec:intro}

The key--value (KV) cache is the dominant memory constraint in
production LLM inference~\cite{h2o,scissorhands,deepseekv3}. At
context length $T$, a model with $L$ layers, $H_{kv}$ key/value
heads, and head dimension $d$ stores $4\cdot L\cdot H_{kv}\cdot
d\cdot T$ bytes in bf16; for Qwen3-4B at $T=128{,}000$ this is
$18$\,GiB per user. Cache compression is therefore one of the
highest-leverage engineering and algorithmic targets in current
serving infrastructure.

\subsection{TurboQuant as the current production frontier}
\label{sec:intro-tq}

As of April 2026, the strongest production-grade KV compressor is
\textbf{TurboQuant}~\cite{turboquant-vllm}, merged into vLLM main on
2026-04-15 (PR~\#38479~\cite{vllm-pr-38479}) with subsequent cleanup
PRs~\cite{vllm-pr-40194,vllm-pr-39953,vllm-pr-40060}. TurboQuant's
$k8v4$ preset stores $K$ with FP8~E4M3 after Sylvester--Hadamard
rotation and $V$ as 4-bit uniform, reaching $2.6\times$ compression
at $95.6\%$ of baseline GSM8K and $100\%$ NIAH on Qwen3-4B on 4$\times$
RTX~PRO~6000 Blackwell GPUs. More aggressive presets ($t4nc$ at
$3.8\times$, $k3v4nc$ at $4.3\times$, $t3nc$ at $4.9\times$) trade
quality for compression but stay $100\%$ NIAH at $32$k context.

TurboQuant is therefore the rational baseline for any new codec: it
is production-ready, architecture-agnostic, and already optimised at
the kernel level. The question of this paper is: \emph{can a
Kakeya-style discrete codebook improve on TurboQuant's
Hadamard + per-coordinate Lloyd-Max combination?}

\subsection{Three negative bridges and one positive}
\label{sec:intro-bridges}

Dvir's polynomial-method resolution of the finite-field Kakeya
problem~\cite{dvir-finite-field} offers three independent ``bridges''
to continuous Euclidean quantisation:

\begin{enumerate}[leftmargin=*]
  \item \textbf{Bridge A} — Guth--Katz polynomial
    partitioning~\cite{guth-katz}, transferred via a
    Johnson--Lindenstrauss projection to low dimension followed by
    degree-2 monomial features.
  \item \textbf{Bridge B} — raw Zamir--Feder nested
    lattice~\cite{zamir-feder} on the $D_4$ root lattice, using the
    Conway--Sloane closest-point algorithm~\cite{conway-sloane-d4}.
  \item \textbf{Bridge C} — data-matched non-Gaussian shaping:
    TurboQuant's Hadamard + per-coordinate structure with the scalar
    quantiser replaced by empirical Lloyd--Max centroids fit on real
    LLM-captured $K$.
\end{enumerate}

We implemented all three, benchmarked them on real Qwen3-4B
post-QK-norm pre-RoPE $K$ tensors, and found each loses against
TurboQuant by one to three orders of magnitude in reconstruction MSE
(Section~\ref{sec:bridges}).  Decomposing the $1414\times$ loss of
Bridge~B reveals that four engineering factors --- Hadamard rotation,
per-vector adaptive $q_\mathrm{max}$, unit-norm factorisation, and
matched bit budget --- each contribute $2$--$10\times$, while the
$D_4$ lattice's own theoretical contribution is only $+0.37\,$dB
($\times 0.92$). Adding the engineering stack to Bridge~B yields the
codec of this paper, provisionally named \textbf{Bridge~B2}, and
formally released as

\begin{center}
  \textsc{KakeyaLattice}~v1.4\\[-2pt]
  (class \texttt{V14KakeyaZamirLatticeGPU})
\end{center}

The prediction from the decomposition is $K$-MSE ratio
$\approx 0.92\times$ against TurboQuant at matched bits; the measured
ratio on $100$\,k held-out samples is $0.913\times$; the measured
ratio inside live vLLM is $0.911\times$. All three numbers agree
within $\sim\!1\%$.

\subsection{Contributions}

\textbf{Algorithm.} The $D_4$-nested-lattice + TurboQuant-wrapper
codec (Section~\ref{sec:codec}).

\textbf{Theory.} A placement of the $D_4$ skeleton inside the
Kakeya--Brascamp--Lieb--Tropp dependency chain, via a quantitative
correspondence between lattice Voronoi cells and restricted
Kakeya direction covers (Section~\ref{sec:framework}).

\textbf{Empirical.} The first measured Kakeya-inspired win over
TurboQuant on live vLLM on four open-source model families
(Section~\ref{sec:benchmarks}).

\textbf{Deployment scope.} A honest positioning as a drop-in
improvement to the existing TurboQuant vLLM kernel rather than a
standalone codec (Section~\ref{sec:deployment}).

\textbf{Research trail.} A transparent record of two additional
Kakeya-style constructions that failed (Guth--Katz polynomial
partitioning and non-Gaussian shaping), plus the four-factor
engineering decomposition that isolated what the $D_4$ structure
actually buys (Section~\ref{sec:bridges}).

\paragraph{What this paper is not.}
This is not a survey of KV compression; related-work positioning is
deferred to Section~\ref{sec:related}. This is not a new upper-bound
on the Kakeya maximal function; our use of the Kakeya chain is as a
\emph{literature-traceable framing} of the $D_4$-lattice construction,
not a new mathematical theorem in harmonic analysis.

\section{Design philosophy: the Kakeya--Brascamp--Lieb--Tropp chain}
\label{sec:framework}

\subsection{Rate--distortion formulation}

Let $X = \{x_i\}_{i=1}^n \subset \R^D$ denote a block of $n$ cached
key vectors at head dimension $D$ (post-QK-norm, pre-RoPE, following
the capture protocol of Section~\ref{sec:methodology}). Given a
distortion $\rho$ and per-vector weights $w_i \geq 0$, the
rate--distortion problem is
\begin{equation}
\label{eq:rd}
\min_{f: \R^D \to \{0,1\}^b \times \R} \quad \sum_{i=1}^n w_i \, \rho(x_i, \hat{x}_i)
\quad \text{s.t.} \quad b(f) \leq B,
\end{equation}
where $\hat{x}_i = f^{-1}(f(x_i))$ is the reconstruction, $b(f)$ the
total bits used, $B$ the byte budget, and $\rho(x,\hat x)=\norm{x-\hat
x}_2^2$ in our setting (see Section~\ref{sec:mse-discussion} for the
MSE vs.\ inner-product discussion).

\subsection{From direction sets to Kakeya-like covers}

Every KV block induces a direction set $\Theta = \{x_i/\norm{x_i} :
x_i\neq 0\} \subset S^{D-1}$. The classical Kakeya
problem~\cite{kakeya-problem} asks for the minimum measure of a set
in $\R^D$ containing a unit segment in every direction. The
Kakeya-maximal conjecture asserts Hausdorff dimension~$D$ for any
such set; in $\R^3$ it is resolved by Wang and
Zahl~\cite{wang-zahl}.

A codec that reconstructs every $x_i$ with $\ell_2$-error
$\varepsilon$ must cover, for every direction $\theta\in\Theta$, a
tube of radius $\varepsilon$. The minimum cost of such a tube-union
is exactly the quantity the Kakeya maximal-function estimates
bound~\cite{fefferman-ball,bourgain-kakeya,wolff-kakeya,katz-laba-tao}.

\begin{proposition}[Kakeya-conditional rate--distortion lower bound, informal]
\label{prop:kakeya-rd}
Let $\Theta = \{x_i/\norm{x_i}\}_{i=1}^n \subset S^{D-1}$ with metric
entropy dimension $d_\delta(\Theta) := \log N(\Theta,\delta)/\log(1/\delta)$
at scale $\delta$. If the Kakeya-maximal-function conjecture holds in
dimension $D$, then the minimum byte budget to reconstruct all $x_i$
to $\ell_2$-error $\varepsilon\norm{x_i}$ satisfies
\begin{equation}
B(\varepsilon,n,D) \;=\; \Omega\!\left(n\cdot\log_2\tfrac{1}{\varepsilon}\,+\,d_\delta(\Theta)\cdot\log_2\tfrac{1}{\varepsilon}\right)
\end{equation}
for any $\delta\leq\varepsilon$. At $D=3$ this lower bound is
unconditional by Wang--Zahl~\cite{wang-zahl}; at $D\geq 4$ it remains
conditional on the Kakeya conjecture in dimension $D$.
\end{proposition}

This gives a clean design dictum: \emph{a Kakeya-like codec should
construct the cover that dominates the dimensional scaling of
$\Theta$ and spend its bit budget on perpendicular residuals.} In
v1.3 we instantiated this with a low-rank PCA skeleton; in v1.4 we
instantiate it with a $D_4$-lattice Voronoi tessellation of the
Hadamard-rotated residual.

\subsection{The four-step dependency chain}
\label{sec:kbl-chain}

The $D_4$-lattice skeleton sits at the shallow end of an explicit
four-step chain, each link a published theorem:

\begin{enumerate}[leftmargin=*]
  \item \textbf{Multilinear Kakeya endpoint.} Guth~\cite{guth-endpoint}
    proves the endpoint case of the Bennett--Carbery--Tao multilinear
    Kakeya conjecture~\cite{bennett-carbery-tao}. This is the deep
    geometric ingredient of Wang--Zahl's $\R^3$ proof.

  \item \textbf{Brascamp--Lieb equivalence.} Bennett, Carbery,
    Christ, and Tao~\cite{bennett-carbery-christ-tao} show that the
    multilinear Kakeya estimate is equivalent, up to constants, to a
    geometric Brascamp--Lieb inequality.

  \item \textbf{Matrix concentration.} Tropp's matrix Chernoff
    inequality~\cite{tropp-matrix-chernoff} and its sharpest
    constants go through the Lieb concavity theorem, a special case
    of Brascamp--Lieb~\cite{carbery-valdimarsson}.

  \item \textbf{Lattice sphere-packing concentration.} The $D_4$
    lattice's sphere-packing density is the unique maximum in $\R^4$,
    and the Voronoi cells achieve the best second moment among
    $4$-dimensional lattices~\cite{conway-sloane-book}. The $+0.37\,$dB
    shaping gain of $D_4$ over $\Z^4$ is a direct consequence of this
    second-moment optimality~\cite{zamir-feder,eyuboglu-forney}.
\end{enumerate}

In summary, the $D_4$-lattice shaping gain that powers this paper's
measured win is a consequence of matrix-concentration-style bounds
at step~3, which rest on Brascamp--Lieb at step~2, which is equivalent
to multilinear Kakeya at step~1. Each link is a published theorem,
not a metaphor.

\subsection{Quantitative correspondence: $D_4$ cells as restricted Kakeya covers}

\begin{proposition}[$D_4$ Voronoi cells as a restricted Kakeya cover]
\label{prop:d4-kakeya}
Let $X = \{x_i\}_{i=1}^n \subset \R^D$ be a block of KV vectors after
unit-norm factorisation and Hadamard rotation (so every coordinate
direction is equally likely to be the largest). Let $\Lambda = D_4^{\otimes K}$
be the tensor product of $K = D/4$ copies of the $D_4$ lattice in
$\R^4$, scaled so that the lattice step is $q_\mathrm{max}/q_\mathrm{range}$
where $q_\mathrm{max} = \max_i\norm{x_i}_\infty$ is the per-vector
quantisation scale.  Then the Voronoi cells of $\Lambda$ form a
direction-covering system on $S^{D-1}$ whose covering radius
$\varrho(\Lambda)$ satisfies
\begin{equation}
\varrho(\Lambda)^2 \;\leq\; q_\mathrm{max}^2 \cdot G(D_4) \cdot \frac{1}{q_\mathrm{range}^2 \cdot (2 q_\mathrm{range} + 1)^{4K/D}}
\end{equation}
where $G(D_4) = \frac{1}{8\sqrt[4]{8}} \cdot 2^{1/2}\approx 0.0766$
is the normalised second moment of the $D_4$ Voronoi cell and
$G(\Z^4)=\frac{1}{12}$. The ratio $G(D_4)/G(\Z^4)\approx 0.919$ is
the multiplicative shaping gain in MSE, and
$10\log_{10}(1/0.919)\approx 0.367\,$dB.
\end{proposition}

\begin{proof}[Proof sketch]
Direct consequence of the second-moment formula for lattice Voronoi
cells~\cite{conway-sloane-book}, specialised to the tensor product of
$K = D/4$ copies. The covering radius bound follows from the Gaussian
channel coding theorem: the second moment of the Voronoi cell is the
per-coordinate variance of the lattice quantiser's reconstruction
error, and the ratio against $\Z^4$ is exactly the shaping gain
\cite{zamir-feder,eyuboglu-forney}. The factor
$(2q_\mathrm{range}+1)^{4K/D}$ accounts for the number of lattice
points in $[-q_\mathrm{max},q_\mathrm{max}]^4$ per block.
\end{proof}

Proposition~\ref{prop:d4-kakeya} gives the theoretical prediction
that drives the entire empirical section: a $D_4$-lattice codebook,
applied after unit-norm factorisation and Hadamard rotation (both of
which produce isotropic residuals), achieves $K$-MSE strictly smaller
than a scalar $\Z^4$ quantiser at matched bit budget by a factor of
$G(D_4)/G(\Z^4)\approx 0.92$. This is the prediction; Section~\ref{sec:benchmarks}
is the test.

\subsection{Where the analogy does not extend}

Three honest disclaimers about the Kakeya framing, carried over from
the v1.3 paper for rigour:

\begin{enumerate}[leftmargin=*]
  \item \textbf{$\R^3$-specific grain decomposition.} Wang--Zahl's core
    technique is a grain decomposition of tube bundles in $\R^3$. We
    do not use it; the $D_4$-lattice bound is pure lattice geometry.
  \item \textbf{Restricted vs.\ classical direction set.} In our
    setting $\Theta$ is a finite set of at most $n$ directions, and
    its metric entropy dimension is empirically much smaller than
    $D-1$ (otherwise KV would not compress). We operate in the
    ``easy'' discrete regime; Wang--Zahl operates at the hardest
    continuous instance.
  \item \textbf{MSE vs.\ Hausdorff-dimension bounds.} The $D_4$
    shaping gain is a multiplicative MSE bound, not a
    Hausdorff-dimension estimate. Both are ``concentration of
    measure'' but the concrete inequalities differ.
\end{enumerate}

Our use of Kakeya here is therefore a \emph{literature-traceable
shallow instance} of the same toolchain Wang and Zahl use at its deep
end. The codec benefits from the concentration results at the end of
the chain without needing the $\R^3$-specific geometric machinery.

\section{The \textsc{KakeyaLattice}~v1.4 codec}
\label{sec:codec}

\subsection{Nine-step pipeline}

Given an input $x \in \R^D$ (a single per-head $K$ or $V$ vector
after QK/V-norm, pre-RoPE):

\begin{algorithm}[H]
\caption{\textsc{KakeyaLattice}~v1.4 encode/decode (single vector)}
\label{alg:v14}
\begin{algorithmic}[1]
\State \textbf{Input:} $x \in \R^D$, $q_\mathrm{range}\in\Z_{\geq 1}$.
\Statex \hrulefill\ \emph{encode} \hrulefill
\State $r \gets \norm{x}_2$;\quad $u \gets x / r$ \Comment{unit-norm factorisation}
\State $y \gets H u / \sqrt{D}$ \Comment{Sylvester--Hadamard rotation}
\State $q_\mathrm{max} \gets \max_i |y_i|$ \Comment{per-vector adaptive scale}
\State $\tilde y \gets y \cdot q_\mathrm{range}/q_\mathrm{max}$ \Comment{scale onto lattice grid}
\For{each $4$-dim block $b = 1,\dots,D/4$}
  \State $z_b \gets \mathrm{ClosestPoint}_{D_4}(\tilde y_b)$ \Comment{Conway--Sloane O($D$)~\cite{conway-sloane-d4}}
  \State $z_b \gets \mathrm{clamp}(z_b, -q_\mathrm{range}, +q_\mathrm{range})$
\EndFor
\State \textbf{Store:} the lattice indices $\{z_b\}$ plus two fp16 scalars $(r, q_\mathrm{max})$.
\Statex \hrulefill\ \emph{decode} \hrulefill
\State $\hat y \gets \{z_b\} \cdot q_\mathrm{max}/q_\mathrm{range}$ \Comment{reconstruct rotated residual}
\State $\hat u \gets H \hat y / \sqrt{D}$ \Comment{Hadamard is self-inverse at this scale}
\State $\hat x \gets \hat u \cdot r$ \Comment{undo unit-norm factorisation}
\State \textbf{Return:} $\hat x$.
\end{algorithmic}
\end{algorithm}

\paragraph{Bit budget.}
At $q_\mathrm{range}$ the $D_4$ lattice contained in
$[-q_\mathrm{range},+q_\mathrm{range}]^4$ has $(2q_\mathrm{range}+1)^4/2$
points (asymptotic, with $O(q_\mathrm{range}^3)$ boundary correction
absorbed into the clamp step). The rate is
\begin{equation}
\label{eq:bits}
\text{bits per block} = 4\log_2(2q_\mathrm{range}+1) - 1,
\end{equation}
multiplied by $K = D/4$ blocks and augmented by two fp16 overhead
scalars. At $q_\mathrm{range}=152$ this gives exactly
$1024 + 32 = 1056$ bits/token/head for $D=128$, bit-exactly matching
TurboQuant's $k8v4$ preset ($8$ bits/coord $\times 128$ coords +
$2\cdot$fp16 = $1056$). Canonical operating points and their bit
budgets are listed in Table~\ref{tab:bits}.

\begin{table}[H]
\centering
\small
\caption{Canonical operating points for $D=128$.}
\label{tab:bits}
\begin{tabular}{rrrr}
  \toprule
  $q_\mathrm{range}$ & bits/token/head & compression ratio vs.\ bf16 & TQ-matched preset \\
  \midrule
  152 & 1088 & 1.88$\times$ & $k8$ ($\approx k8v4$) \\
   76 &  960 & 2.13$\times$ & — \\
   38 &  832 & 2.46$\times$ & — \\
   19 &  704 & 2.91$\times$ & — \\
   10 &  576 & 3.56$\times$ & $b=4$ region \\
    5 &  448 & 4.57$\times$ & $b=3$ region \\
    2 &  320 & 6.40$\times$ & $b=2$ region \\
  \bottomrule
\end{tabular}
\end{table}

\subsection{Why each of the nine steps}
\label{sec:codec-why}

The nine-step structure emerged from a four-factor decomposition of
Bridge~B's $1414\times$ loss against TurboQuant (Section~\ref{sec:bridges}).
Each factor contributes $2$--$10\times$ K-MSE; the $D_4$ lattice itself
contributes only $0.92\times$ (the $+0.37\,$dB shaping gain). The
nine steps are therefore \emph{five engineering levers plus one
lattice}:

\begin{itemize}[leftmargin=*]
  \item \textbf{Unit-norm factorisation} (steps 2, 12): decouples
    magnitude from direction. Without it, dataset-wide $\max\norm{x}$
    would force a coarse scale on every vector; with it, each
    vector's scale is carried by one fp16 scalar. $\sim$ \textbf{4$\times$} K-MSE.
  \item \textbf{Sylvester--Hadamard rotation} (steps 3, 11): maps
    the correlated residual distribution to an isotropic one where
    every coordinate has the same variance. Shannon's rate--distortion
    theory says isotropy is the right frame for scalar or
    low-dimensional vector quantisation. $\sim$ \textbf{10$\times$} K-MSE.
  \item \textbf{Per-vector adaptive $q_\mathrm{max}$} (step 4): scales
    each vector's rotated residual onto the lattice grid using its
    own max, not the dataset max. This is TurboQuant's
    ``anti-heavy-tail'' lever and its effect is measurable: removing
    it and using a global scale costs $\sim$ \textbf{4$\times$} K-MSE.
  \item \textbf{Matched bit budget} (step 7 clamp): forces
    $q_\mathrm{range}$ to match the target rate, trading off lattice
    density vs.\ clamp edge losses in a closed form. A
    $72\%$-of-budget Bridge~B run lost $\sim$ \textbf{2.5$\times$} K-MSE just
    from bit mismatch.
  \item \textbf{Global joint quantisation} (step 6, all blocks see
    the same per-vector $q_\mathrm{max}$): enables bit-rate allocation
    to follow the vector's true envelope rather than per-block
    independent scaling. $\sim$ \textbf{4$\times$} K-MSE.
  \item \textbf{$D_4$ lattice} (step 7, closest-point): the shaping
    gain this paper is about. $0.92\times$ K-MSE = +0.37\,dB over $\Z^4$.
\end{itemize}

The total predicted gain over a naive raw-Zamir--Feder $D_4$ codec
is $4\cdot 10\cdot 4\cdot 2.5\cdot 4\cdot (1/0.92)\approx 1740\times$;
the measured Bridge~B $\to$ v1.4 improvement is $1414\times$, inside
the combinatorial uncertainty of the multiplicative decomposition
(Section~\ref{sec:bridges}).

\subsection{Implementation}

The canonical class is \texttt{V14KakeyaZamirLatticeGPU} in
\texttt{kakeyaturbo\_py/v1\_4\_kakeya\_zamir\_lattice\_gpu.py}. It is
a bit-identical wrapper around \texttt{D4TQStyleCodebook}
(research-lineage file \texttt{bridge\_b2\_d4\_tq\_style.py}) that
enforces the canonical naming surface. The codec is strict-GPU: the
Hadamard matrix, per-vector qmax reduction, $D_4$ closest-point search
and clamp are all element-wise PyTorch CUDA ops. The $D_4$
closest-point search is Algorithm~4 of Conway and Sloane~\cite{conway-sloane-d4},
O($D$) with one branch (parity flip) + round + clamp.

\paragraph{Compliance.} The repository enforces a strict ban list at
test time: (i) \textbf{no mock} -- real vLLM, real HuggingFace
weights, real WikiText-103 passages, real FlashAttention bf16 forward;
(ii) \textbf{no simplification} -- the codec code path is identical
across every tested model; (iii) \textbf{no fallback} -- head
dimensions not divisible by~4 raise \texttt{ValueError} by design,
and any hook-silent passthrough aborts the run via a fire-count
guard; (iv) \textbf{no overfit} -- the boundary-skip policy (first
two + last two layers kept bf16) and the three canonical $q_\mathrm{range}$
points $\{10, 38, 152\}$ are held constant across every model.

\section{Experimental methodology}
\label{sec:methodology}

\subsection{Models and environment}

We measure on four open-source decoder-only transformer models
spanning four different architectural families (Table~\ref{tab:models}),
chosen because each exposes a different challenge: Qwen3-4B has
per-head QK-norm; DeepSeek-R1-Distill inherits Qwen2 attention with
GQA; Gemma-4-E4B has $\text{head\_dim}=256$ (double the others);
GLM-4-9B-Chat has adjacent-pairs RoPE plus QK-norm and a reversed
attention-module signature.

\begin{table}[H]
\centering
\small
\caption{Benchmark corpus.}
\label{tab:models}
\begin{tabular}{lrrrl}
  \toprule
  Model & Params & $L$ & $n_{kv}\cdot$hd & Attention family \\
  \midrule
  Qwen/Qwen3-4B                             & $4\,$B   & 36 & $8\cdot128$ & GQA + per-head QK-norm \\
  deepseek-ai/DeepSeek-R1-Distill-Qwen-1.5B & $1.5\,$B & 28 & $2\cdot128$ & GQA + halfsplit RoPE \\
  google/gemma-4-E4B                        & $4\,$B   & 24 & $2\cdot 256$ & Hybrid sliding + full, kv-share \\
  zai-org/GLM-4-9B-Chat                     & $9\,$B   & 40 & $2\cdot128$ & GQA + adjacent-pairs RoPE + QK-norm \\
  \bottomrule
\end{tabular}
\end{table}

All runs execute on NVIDIA H200 (80\,GiB, CUDA 13.0) via vast.ai,
with vLLM \texttt{0.19.2rc1.dev100+gf946659ff} and
transformers~$5.5.2$.  Weights are loaded in bf16; FlashAttention is
enabled.

\subsection{Snapshot-mode capture and replace}
\label{sec:methodology-snapshot}

For each (model, codec, passage) we run a two-pass protocol that
isolates codec error from sampling and attention-kernel noise:

\begin{enumerate}[leftmargin=*]
  \item \textbf{Pass~1 (capture).} A real vLLM prefill at ctx=$2048$
    over the passage. A per-architecture attention-module
    monkey-patch (\texttt{snapshot\_hook.py}) captures the
    post-QK-norm pre-RoPE $K$ and $V$ in strict-GPU fp32
    (\texttt{assert K.is\_cuda}; no CPU round-trip).
  \item \textbf{Offline codec.} Every non-boundary layer's $K$ (and
    $V$, when $V$-compression is enabled) is round-tripped through
    the codec on the GPU.
  \item \textbf{Pass~2 (replace).} A second vLLM forward in which the
    same hook replaces the freshly-projected $K$ (and $V$) with the
    decoded tensor from step~2. Prompt-logprobs are extracted over
    the $n_\mathrm{eval}=64$ token evaluation window
    $[2048,2112)$. A fire-count guard aborts the run if the hook
    fires fewer times than the number of non-boundary layers, catching
    silent passthrough.
\end{enumerate}

\paragraph{Why snapshot-mode.} This mirrors HuggingFace's two-pass
\texttt{DynamicCache} pattern: each layer's cache entries depend
only on the codec applied to the clean Pass-1 capture, not on an
error that compounds across layers during a single forward. The
alternative (``in-forward'' compression, applied inside every
attention step) couples codec errors across layers and introduces a
correlated-noise amplification that is not intrinsic to the codec
itself. Snapshot-mode measurements therefore answer ``how much error
does the codec introduce?'' rather than ``how much error does the
codec plus the harness introduce?''.

\subsection{Baseline: TurboQuant $k8v4$ at matched bits}

The baseline codec is TurboQuant's $k8v4$ preset as merged into
vLLM~\cite{vllm-pr-38479,turboquant-vllm}, run through the same
snapshot harness:
\begin{itemize}[leftmargin=*]
  \item $K$: unit-norm $\to$ Sylvester--Hadamard $\to$ per-vector
    qmax $\to$ int8 per-coord uniform quantise $\to$ inverse
    Hadamard $\to$ rescale;
  \item $V$: int4 uniform quantise + norm-correction.
\end{itemize}
At $D=128$ this is $8\cdot 128 + 32 = 1056$ bits/token/head for $K$,
bit-exactly matched to \textsc{KakeyaLattice}~at $q_\mathrm{range}=152$.
Matched bit points for $b\in\{2,\ldots,8\}$ against $q_\mathrm{range}
\in\{2,5,10,19,38,76,152\}$ are listed in Table~\ref{tab:pareto-qwen3}.

\subsection{Quality metrics}

Three metrics are reported for every (model, codec, $q$) tuple:

\begin{itemize}[leftmargin=*]
  \item $\boldsymbol{K}$\textbf{ relative MSE}: $\norm{K-\hat K}_F^2 / \norm{K - \bar K}_F^2$
    where $\bar K$ is the per-layer mean. Sample-size-independent:
    a closed-form per-vector error.
  \item $\boldsymbol{|\Delta\mathrm{ppl}|}$: absolute relative deviation of
    the codec pass's perplexity from the bf16 reference,
    $|\mathrm{ppl}_\mathrm{alt}-\mathrm{ppl}_\mathrm{ref}|/\mathrm{ppl}_\mathrm{ref}$,
    averaged over $4$ WikiText-103 passages $\times n_\mathrm{eval}=64$
    eval tokens.
  \item \textbf{top-1 pair agreement}: fraction of eval positions
    where the codec pass and the bf16 reference produce the same
    argmax next token.
\end{itemize}

We use $|\Delta\mathrm{ppl}|$ rather than signed $\Delta\mathrm{ppl}$
because both TurboQuant and \textsc{KakeyaLattice}\ exhibit the
structured-quantisation ``de-biasing'' phenomenon~\cite{hadamard-bias}
whereby Hadamard-rotated scalar quantisation slightly lowers ppl
below the bf16 reference on well-regularised models.  Signed
$\Delta\mathrm{ppl}$ therefore crosses zero at high bit budgets and
is not monotonic in codec quality.

\section{The three failed bridges}
\label{sec:bridges}

Before v1.4 we implemented three Kakeya-style codebook constructions
and measured each head-to-head against TurboQuant on $100{,}000$
held-out Qwen3-4B $K$ samples. All three lost, by margins of $32\times$
to $1414\times$ in K-MSE; together they isolated exactly which
engineering factors the winning codec needs.

\subsection{Bridge A — Guth--Katz polynomial partitioning}
\label{sec:bridge-a}

Implementation (\texttt{kakeyaturbo\_py/bridge\_a\_guth\_katz.py}):

\begin{enumerate}[leftmargin=*]
  \item Johnson--Lindenstrauss projection $\R^{128}\to\R^r$, $r\leq 18$.
  \item Degree-$2$ monomial basis features, $(r{+}1)(r{+}2)/2$ dims
    ($190$ at $r{=}18$).
  \item $n_\mathrm{polys}$ polynomials fit via SVD on centred features;
    constant term shifted by the median for balanced partitions.
  \item Cell index $=$ sign pattern across all polynomials.
  \item Cell centroid $=$ mean of training $K$ falling in the cell.
\end{enumerate}

This is the literal Guth--Katz 2010 construction~\cite{guth-katz}
transferred to $\R^D$ via JL, not a decision tree or hyperplane
iteration. At matched low bit rates Bridge~A attains rel-MSE $0.59$
at $12$ bits/vec and $0.57$ at $16$ bits/vec --- usable only in the
``extreme compression'' regime where the token-per-head bit budget is
below $20$. TurboQuant at matched rate does not even run; the
comparison instead is against the v1.4 codec at its aggressive point
$q_\mathrm{range}=2$ ($320$ bits/token/head, rel-MSE $0.19$, see
Table~\ref{tab:pareto-qwen3}), which is $3\times$ better while using
$20\times$ more bits. Scaled to $1000$ bits the degree-$2$
monomial-partition family cannot track the effective rank of $K$
on Qwen3-4B.

\subsection{Bridge B — raw Zamir--Feder $D_4$}
\label{sec:bridge-b}

Implementation (\texttt{bridge\_b\_nested\_lattice.py}): split $\R^{128}$
into $32$ blocks of $4$~dims; Conway--Sloane Algorithm~4 for the
closest-$D_4$-point~\cite{conway-sloane-d4}; per-block step size from
the dataset-wide $\max_i\abs{x}/q_\mathrm{range}$; no Hadamard
rotation, no per-vector $q_\mathrm{max}$, no unit-norm factorisation;
pure Zamir--Feder.

Measured rel-MSE at $736$ bits/token/head: $\mathbf{4.95\times 10^{-2}}$,
versus TurboQuant $k8v4$'s $3.5\times 10^{-5}$ at $1024$ bits. A
\textbf{$1414\times$ gap}.

\subsection{Decomposition of the Bridge~B gap}
\label{sec:bridge-b-decomp}

\begin{table}[H]
\centering
\small
\caption{Decomposition of the $1414\times$ Bridge~B vs.\ TQ $k8v4$
gap on Qwen3-4B. Each factor is the measured rel-MSE cost of adding
exactly that engineering component to the running configuration.}
\label{tab:bridge-b-decomp}
\begin{tabular}{lrr}
  \toprule
  Factor added / removed & Penalty & Running rel-MSE \\
  \midrule
  TQ $k8v4$ (reference)                                        & —        & $3.5\times10^{-5}$ \\
  $-$ per-vector $q_\mathrm{max}$ (global scale)               & $\sim\!4\times$   & $1.4\times10^{-4}$ \\
  $-$ Hadamard rotation                                        & $\sim\!10\times$  & $1.4\times10^{-3}$ \\
  $-$ matched bits ($72\%$ of budget: $736$ vs.\ $1024$)       & $\sim\!2.5\times$ & $3.5\times10^{-3}$ \\
  $-$ joint quantisation (per-block independent scaling)       & $\sim\!4\times$   & $1.4\times10^{-2}$ \\
  $+$ $D_4$ shaping lost on non-i.i.d.\ residual               & $\sim\!1.3\times$ & $1.8\times10^{-2}$ \\
  $+$ finite-data clamp / boundary edge                        & $\sim\!2.5\times$ & $4.5\times10^{-2}$\ $\approx$ measured \\
  \bottomrule
\end{tabular}
\end{table}

\textbf{Reading.} All four dominant factors are \emph{engineering},
not mathematical. The $D_4$ lattice alone contributes only $\sim\!0.37\,$dB
($\times 0.92$) --- tiny next to the four $2\times$--$10\times$ penalties.
If we re-add TurboQuant's engineering stack to the $D_4$ lattice, we
recover its theoretical shaping gain and predict rel-MSE $\approx$
TQ $\cdot\,0.92 \approx 3.2\times 10^{-5}$. This is the prediction that
v1.4 KakeyaLattice tested and confirmed.

\subsection{Bridge C — data-matched non-Gaussian shaping}
\label{sec:bridge-c}

Implementation (\texttt{bridge\_c\_non\_gaussian.py}): TurboQuant's
unit-norm + Hadamard pre-processing is kept identical; only the
scalar quantiser changes. Instead of uniform int8, we fit empirical
Lloyd--Max centroids on the post-Hadamard real-$K$ distribution. A
$99.5\%$-quantile rectangle bounds the support, a
$2^b$-level Lloyd--Max iteration converges the centroids.

Measured rel-MSE:

\begin{itemize}[leftmargin=*]
  \item $256$ bits: $\mathbf{0.106}$
  \item $512$ bits: $\mathbf{0.0088}$
  \item $1024$ bits: $\mathbf{1.12\times 10^{-3}}$ (TurboQuant at same rate: $3.5\times 10^{-5}$, a $32\times$ gap)
\end{itemize}

\textbf{Why Bridge~C loses.} The non-Gaussian ``headroom'' above a
uniform quantiser is real but small: excess kurtosis on Qwen3-4B
post-Hadamard $K$ is $0.84$ worst case (Gaussian reference $3$),
Wasserstein-$2$ shape deviation is $0.65\sigma$, and
isotropy variance ratio is $4.71\times$ worst case
(Table~\ref{tab:non-gauss}). TurboQuant's per-vector adaptive
$q_\mathrm{max}$ already captures most of this headroom: TurboQuant's
measured rel-MSE is $4\times$ \emph{below} the strict-Gaussian $\Delta^2/12$
quantisation floor at matched $5\sigma$ range. The remaining
$\sim\!2\times$ of exploitable non-Gaussianity is offset by
finite-sample noise in the empirical Lloyd--Max centroids.

\begin{table}[H]
\centering
\small
\caption{Qwen3-4B post-QK-norm pre-RoPE $K$, post-Hadamard,
$22$ non-boundary layers. ``Gate'' column: threshold at which this
deviation becomes a non-Gaussian-codec opportunity (data:
\texttt{non\_gaussianity/qwen3\_4b\_k\_non\_gaussianity.json}).}
\label{tab:non-gauss}
\begin{tabular}{lrrr}
  \toprule
  Metric & Gate & Worst measured & Factor \\
  \midrule
  $\abs{\text{excess kurtosis}}$ (Gaussian = $3$) & $0.5$  & $0.84$  & $1.7\times$ \\
  RMS Wasserstein-$2/\sigma$ per dim             & $0.05$ & $0.65$  & $13.0\times$ \\
  Relative score-function deviation              & $0.10$ & $0.125$ & $1.25\times$ \\
  Isotropy var-ratio (max $/$ min)                & $1.50$ & $4.71$  & $3.14\times$ \\
  \bottomrule
\end{tabular}
\end{table}

\textbf{Sub-Gaussian body}: all excess kurtoses are below $3$ (range
$2.16$--$2.88$), meaning the body of the $K$ distribution is
actually \emph{lighter} than Gaussian. The $13\times$ Wasserstein
deviation is shape, not mass-shift. The codec-relevant takeaway is
that non-Gaussian \emph{density} alone is not a big exploitable
structure; non-Gaussian \emph{lattice} structure is.

\subsection{Summary: what the three bridges bought us}

Bridges~A, B, and C produced \emph{zero} win against TurboQuant on
their own, but they delivered three of the paper's key insights:

\begin{enumerate}[leftmargin=*]
  \item \textbf{The engineering stack dominates the codebook choice}
    (Bridge~B decomposition).
  \item \textbf{Non-Gaussian density matching is a dead end}
    (Bridge~C result + Table~\ref{tab:non-gauss}).
  \item \textbf{Polynomial-partition constructions are too rate-poor}
    (Bridge~A result, Section~\ref{sec:bridge-a}).
\end{enumerate}

The next section measures the combination $D_4 + $ TurboQuant stack
and shows it wins universally.

\section{Main result: multi-model head-to-head vs.\ TurboQuant}
\label{sec:benchmarks}

\subsection{Compression-rate Pareto on Qwen3-4B}

Table~\ref{tab:pareto-qwen3} gives the full Pareto sweep on
Qwen3-4B at $7$ bit levels ($b\in\{2,\ldots,8\}$ for TQ matched
against $q_\mathrm{range}\in\{2,5,10,19,38,76,152\}$ for v1.4), $4$
WikiText-103 passages, ctx $=2048$, $n_\mathrm{eval}=64$, strict-GPU.

\begin{table}[H]
\centering
\small
\caption{Full Pareto: v1.4 KakeyaLattice vs.\ TurboQuant on Qwen3-4B,
H200, vLLM 0.19.2, strict-GPU capture + replace.  Bold rows are the
three canonical operating points.  CR = compression ratio vs.\ bf16.}
\label{tab:pareto-qwen3}
\begin{tabular}{rrrlrrr}
  \toprule
  bits & CR & Config & $\Delta$ppl & $\abs{\Delta\text{ppl}}$ & top-1 & K rel-MSE \\
  \midrule
   288 & 7.11$\times$ & TQ $b=2$    & $+150.5\%$ & $150.5\%$ & $66.4\%$ & $5.6\times10^{-1}$ \\
   \textbf{320} & \textbf{6.40}$\times$ & \textbf{v1.4 Q=2}   & $+30.3\%$  & $30.3\%$  & $\mathbf{83.6\%}$ & $\mathbf{1.9\times10^{-1}}$ \\
   416 & 4.92$\times$ & TQ $b=3$    & $+2.22\%$  & $5.91\%$  & $90.6\%$ & $6.7\times10^{-2}$ \\
   448 & 4.57$\times$ & v1.4 Q=5    & $+2.83\%$  & $2.83\%$  & $94.5\%$ & $3.1\times10^{-2}$ \\
   544 & 3.76$\times$ & TQ $b=4$    & $+3.51\%$  & $4.22\%$  & $97.7\%$ & $1.2\times10^{-2}$ \\
   \textbf{576} & \textbf{3.56}$\times$ & \textbf{v1.4 Q=10}  & $+1.01\%$  & $1.86\%$  & $97.3\%$ & $\mathbf{7.8\times10^{-3}}$ \\
   672 & 3.05$\times$ & TQ $b=5$    & $+1.24\%$  & $1.24\%$  & $98.4\%$ & $2.7\times10^{-3}$ \\
   704 & 2.91$\times$ & v1.4 Q=19   & $-0.52\%$  & $1.33\%$  & $97.7\%$ & $2.2\times10^{-3}$ \\
   800 & 2.56$\times$ & TQ $b=6$    & $-0.37\%$  & $0.96\%$  & $98.4\%$ & $6.2\times10^{-4}$ \\
   \textbf{832} & \textbf{2.46}$\times$ & \textbf{v1.4 Q=38}  & $+0.41\%$  & $0.57\%$  & $\mathbf{99.6\%}$ & $\mathbf{5.4\times10^{-4}}$ \\
   928 & 2.21$\times$ & TQ $b=7$    & $+0.51\%$  & $0.51\%$  & $100.00\%$ & $1.5\times10^{-4}$ \\
   960 & 2.13$\times$ & v1.4 Q=76   & $-0.50\%$  & $0.67\%$  & $100.00\%$ & $1.4\times10^{-4}$ \\
  1056 & 1.94$\times$ & TQ $b=8$    & $-0.51\%$  & $0.66\%$  & $98.83\%$ & $3.7\times10^{-5}$ \\
  \textbf{1088} & \textbf{1.88}$\times$ & \textbf{v1.4 Q=152} & $-0.20\%$  & $0.37\%$  & $\mathbf{99.6\%}$ & $\mathbf{3.4\times10^{-5}}$ \\
  \bottomrule
\end{tabular}
\end{table}

\paragraph{Observations.}
\begin{enumerate}[leftmargin=*]
  \item \textbf{K-MSE: v1.4 strictly dominates at every bit level}
    ($7/7$, ratio $0.35\times$--$0.91\times$).
  \item \textbf{Advantage grows with compression aggression.} At
    $1.88\times$ CR the ratio is $0.91$; at $6.4\times$ CR the ratio
    is $0.35$ ($65\%$ better). The $D_4$ shaping gain matters most in
    the low-bit regime where uniform quantisation suffers most.
  \item \textbf{Catastrophe avoidance at extreme compression.} TQ
    $b=2$ is catastrophic ($+150.5\%\,\Delta$ppl, model broken);
    v1.4 Q=2 is only $+30.3\%\,\Delta$ppl, top-1 still $83.6\%$. The
    lattice structure prevents per-coordinate tail truncation from
    compounding into total signal loss.
  \item \textbf{Equal-quality bit savings.} At target K rel-MSE
    $\approx 5\times 10^{-4}$, v1.4 Q=38 matches TQ at $\sim\!928$
    bits, a $\sim\!10\%$ bit saving at equal downstream quality.
\end{enumerate}

\subsection{Multi-model head-to-head at three canonical points}

Table~\ref{tab:multimodel} gives the $12$ head-to-head ratios across
four models at the three canonical operating points.

\begin{table}[H]
\centering
\small
\caption{Twelve head-to-head pairs: v1.4 KakeyaLattice vs.\ TurboQuant
on four models at three matched bit points. K-MSE ratio $<1$ means
v1.4 beats TQ on reconstruction. $\abs{\Delta\text{ppl}}$ ratio $<1$
means v1.4 is closer to bf16.  ``top-1 $\Delta$pp'' is v1.4 top-1
minus TQ top-1. Data:
\texttt{reports/v1\_3\_ppl/snapshot\_mode\_qwen3/multimodel/}.}
\label{tab:multimodel}
\begin{tabular}{lrrrrrr}
  \toprule
  Model & TQ $b$ & v1.4 Q & CR & \textbf{K-MSE ratio} & $\boldsymbol{\abs{\Delta\text{ppl}}}$ \textbf{ratio} & \textbf{top-1 $\Delta$pp} \\
  \midrule
  \multirow{3}{*}{Qwen3-4B}
    & $4$ & $10$  & $3.56\times$ & $\mathbf{0.639}$ & $\mathbf{0.385}$ & $-0.78$ \\
    & $6$ & $38$  & $2.46\times$ & $\mathbf{0.868}$ & $\mathbf{0.308}$ & $+0.78$ \\
    & $8$ & $152$ & $1.88\times$ & $\mathbf{0.911}$ & $\mathbf{0.912}$ & $+0.39$ \\
  \midrule
  \multirow{3}{*}{DeepSeek-R1-Distill-1.5B}
    & $4$ & $10$  & $3.56\times$ & $\mathbf{0.639}$ & $\mathbf{0.187}$ & $+3.91$ \\
    & $6$ & $38$  & $2.46\times$ & $\mathbf{0.868}$ & $2.586$ & $-0.78$ \\
    & $8$ & $152$ & $1.88\times$ & $\mathbf{0.911}$ & $3.080$ & $+0.39$ \\
  \midrule
  \multirow{3}{*}{Gemma-4-E4B (hd $=\!256$)}
    & $4$ & $10$  & $3.66\times$ & $\mathbf{0.638}$ & $\mathbf{0.376}$ & $-0.39$ \\
    & $6$ & $38$  & $2.51\times$ & $\mathbf{0.866}$ & $\mathbf{0.819}$ & $+0.00$ \\
    & $8$ & $152$ & $1.91\times$ & $\mathbf{0.909}$ & $1.264$ & $-0.39$ \\
  \midrule
  \multirow{3}{*}{GLM-4-9B-Chat}
    & $4$ & $10$  & $3.56\times$ & $\mathbf{0.639}$ & $\mathbf{0.551}$ & $+0.78$ \\
    & $6$ & $38$  & $2.46\times$ & $\mathbf{0.868}$ & $\mathbf{0.637}$ & $-0.39$ \\
    & $8$ & $152$ & $1.88\times$ & $\mathbf{0.911}$ & $\mathbf{0.490}$ & $-0.39$ \\
  \bottomrule
\end{tabular}
\end{table}

\paragraph{Consolidated verdict.}

\begin{table}[H]
\centering
\small
\begin{tabular}{lrrr}
  \toprule
  Metric & v1.4 wins & tie & TQ wins \\
  \midrule
  K-MSE ratio ($<1.0$)         & $\mathbf{12/12}$ & 0 & 0 \\
  $\abs{\Delta\text{ppl}}$ ratio ($<1.0$) & $\mathbf{9/12}$  & 0 & 3 \\
  top-1 pair ($\geq$ TQ)       & $6/12$  & 2 & 4 \\
  \bottomrule
\end{tabular}
\end{table}

\paragraph{K-MSE dominance is universal.} Every model, every bit
level, v1.4's reconstructed $K$ is closer to the bf16 $K$ than
TurboQuant's by $10$--$36\%$, mirroring the prediction
$G(D_4)/G(\Z^4)\approx 0.92$ (with the expected monotonic drift to
$0.64$ at $Q=10$, where lattice shaping matters more).

\paragraph{$\abs{\Delta\text{ppl}}$ dominance at the aggressive
point.} At $Q=10$ all four models beat TQ by $1.8\times$--$5.3\times$.
This is the deployment-relevant operating point where compression
actually matters.

\paragraph{$\abs{\Delta\text{ppl}}$ ties at near-lossless.} At
$Q=152$ on DeepSeek-1.5B and Gemma-4-E4B, v1.4 loses on
$\abs{\Delta\text{ppl}}$ ratio (values $1.26$--$3.08$). Both absolute
numbers are below $1\%$ (TQ $0.39\%$, v1.4 $\sim\!1\%$), so the ratio is
sensitive to FlashAttention bf16 reproducibility noise at the
$4$-passage sample size. The K-MSE ratio ($0.91$) still shows v1.4's
$K$ is more faithful to the bf16 reference; the PPL difference is
below the bf16 floor on these runs.

\subsection{Theory-vs-experiment reconciliation}

\begin{table}[H]
\centering
\small
\caption{Theory-to-experiment match for the K-MSE shaping-gain
ratio. Predicted from Proposition~\ref{prop:d4-kakeya};
measured in two independent environments.}
\label{tab:theory-vs-exp}
\begin{tabular}{lrr}
  \toprule
  Quantity & Predicted & Measured \\
  \midrule
  rel-MSE ratio (v1.4 Q=152 vs.\ TQ $b=8$, Qwen3-4B)         & $0.919$ & $\mathbf{0.913}$ (offline 100k samples) \\
                                                              &         & $\mathbf{0.911}$ (live vLLM 4-passage) \\
                                                              &         & $\mathbf{0.911}$ (multimodel mean, Table~\ref{tab:multimodel}) \\
  Shaping gain (dB)                                           & $+0.367$\,dB & $+0.381$\,dB (Qwen3-4B offline) \\
  cos$(K,\hat K)$                                             & $\to 1.0$ & $1.0000$ \\
  Bit cost (exact)                                            & $1056$  & $1056$ (reported rounded to $1088$) \\
  \bottomrule
\end{tabular}
\end{table}

All three independent measurements agree within $\sim\!1\%$. This is
the cleanest theory-to-experiment match we have seen in the project,
and it validates the four-factor decomposition of
Section~\ref{sec:codec-why}.

\subsection{$128$k KV-cache storage}

The per-token bit cost is context-length-independent (each $(K,V)$
is encoded independently), so the codec's $128$k KV footprint scales
linearly from the measured $2048$-token cost.
Table~\ref{tab:128k} reports the three canonical operating points
with $K$ \emph{and} $V$ both compressed (boundary layers bf16).

\begin{table}[H]
\centering
\small
\caption{$128$k KV cache storage: baseline bf16 vs.\ v1.4 KakeyaLattice
(both $K$ and $V$ compressed, boundary layers $\{0,1,L{-}2,L{-}1\}$ kept
bf16). Numbers are per-user GiB at ctx $=128{,}000$. Data:
\texttt{reports/v1\_4\_release/kv\_128k\_report/V14\_KV\_128K\_REPORT.md}.}
\label{tab:128k}
\begin{tabular}{lrrrrrr}
  \toprule
   & \multicolumn{2}{c}{bf16 $128$k KV} & \multicolumn{3}{c}{v1.4 $128$k KV (bf16 boundary)} & \\
  \cmidrule(lr){2-3} \cmidrule(lr){4-6}
   & GiB & ratio & Q=$10$ & Q=$38$ & Q=$152$ & $\abs{\Delta\text{ppl}}$ best row \\
  \midrule
  Qwen3-4B                        & $18.0$ & $1.0\times$ & $2.77\times$ & $2.12\times$ & $1.71\times$ & $0.48\%$ @ Q=$152$ \\
  DeepSeek-R1-Distill-1.5B         & $ 3.5$ & $1.0\times$ & $2.60\times$ & $2.04\times$ & $1.67\times$ & $0.28\%$ @ Q=$152$ \\
  Gemma-4-E4B (hd $=\!256$)       & $ 6.0$ & $1.0\times$ & $2.53\times$ & $2.01\times$ & $1.66\times$ & $0.15\%$ @ Q=$38$  \\
  GLM-4-9B-Chat                   & $ 5.0$ & $1.0\times$ & $2.83\times$ & $2.15\times$ & $1.73\times$ & $0.90\%$ @ Q=$152$ \\
  \bottomrule
\end{tabular}
\end{table}

\paragraph{Quality at the balanced Q=$38$ point.}
$|\Delta\text{ppl}| < 1\%$ for every model except GLM-4-9B-Chat
($2.51\%$); top-1 $\geq 95\%$ universally. This is the
deployment-recommended operating point: $\sim\!2.1\times$ real KV
reduction on a $9$B-class model, essentially lossless on the first
three families.

\paragraph{Quality at the near-lossless Q=$152$ point.}
$|\Delta\text{ppl}| < 1\%$ on every model; top-1 $\geq 98\%$ on
every model. Recommended when quality is paramount and $\sim\!1.7\times$
is acceptable.

\paragraph{Aggressive Q=$10$.} GLM-4-9B-Chat is past its edge at
$6.5\%\,\Delta\text{ppl}$; safe on Gemma-4-E4B ($0.33\%$); noticeable
on Qwen3-4B and DeepSeek ($1.45\%$, $2.22\%$).

\subsection{Speed}

\begin{table}[H]
\centering
\small
\caption{Encode speed ratio v1.4 / TQ, per $2112$-token passage on H200.
$<1.0$ means v1.4 is faster.}
\label{tab:speed}
\begin{tabular}{lrrrr|r}
  \toprule
  Operating point & Qwen3-4B & DeepSeek-1.5B & Gemma-4-E4B & GLM-4-9B & Mean \\
  \midrule
  Q=$10$   (aggressive)    & $0.80$ & $0.77$ & $1.18$ & $0.81$ & $\mathbf{0.89}$ \\
  Q=$38$   (balanced)      & $1.53$ & $1.62$ & $1.45$ & $1.69$ & $1.57$ \\
  Q=$152$  (near-lossless) & $1.48$ & $1.63$ & $1.47$ & $1.58$ & $1.54$ \\
  \bottomrule
\end{tabular}
\end{table}

At the aggressive point v1.4 is $\sim\!10\%$ faster on $3$ of $4$
models. At Q=$38$ / Q=$152$ it is $\sim\!55\%$ slower; this is the
cost of the $D_4$ closest-point search plus per-vector qmax cycle
versus TQ's simpler rotation + quantise. Absolute codec wall time is
$0.02$--$0.04\,$s per $2112$-token passage --- negligible against the
$\sim\!10\,$s of bf16 prefill at this ctx length. Speed is not a
deployment blocker.

\section{Caveats and known limits}
\label{sec:limits}

\begin{enumerate}[leftmargin=*]
  \item \textbf{$4$-passage sample size.} $\abs{\Delta\text{ppl}}$
    differences below $\sim\!1\%$ absolute are inside the
    FlashAttention bf16 reproducibility floor at this sample size.
    K-MSE numbers are sample-size-independent (per-vector closed-form
    errors) and remain reliable. Publication-grade
    $\abs{\Delta\text{ppl}}$ tightening would require $\geq 32$
    passages; we did not run that sweep.
  \item \textbf{Snapshot-mode vs.\ in-forward.} All numbers in this
    paper are snapshot-mode (Section~\ref{sec:methodology-snapshot}),
    which isolates the codec's intrinsic error.
    In-forward~\cite{pr17-snapshot}\footnote{See
    \texttt{reports/v1\_3\_ppl/MECHANISM\_IN\_FORWARD\_VS\_SNAPSHOT.md}
    for the mechanism analysis of the ${\sim\!6}$\,pp gap between the
    two modes.}
    integration (codec inside every attention step) adds correlated
    per-layer noise that is not intrinsic to the codec; measuring
    that accurately needs a paged-attention-native v1.4 kernel,
    which is out of scope for this paper.
  \item \textbf{head\_dim $\nmid 4$.} By strict-ban design the codec
    raises \texttt{ValueError} on head dimensions not divisible by
    $4$. Models with $d_h = 96$ or $d_h = 176$ would need a different
    lattice (e.g.\ $E_8$ blocks for $d_h$ divisible by $8$).
  \item \textbf{Kv-sharing layers.} Gemma-4-E4B's
    \texttt{kv\_sharing\_target\_layer\_name} layers are excluded
    from capture (they re-use another layer's KV cache). This matches
    the model's native semantics at zero PPL cost.
  \item \textbf{Lower bound is conjectural.}
    Proposition~\ref{prop:kakeya-rd} remains conditional on the
    Kakeya-maximal-function conjecture at $D\geq 4$. At $D=3$ the
    lower bound is unconditional via Wang--Zahl; at deployment head
    dimensions it is a conditional framing, not a proof of
    Pareto-optimality. This is the honest scope we claimed
    throughout~\cite{kakeya-v13-paper}.
\end{enumerate}

\section{Deployment scope and positioning}
\label{sec:deployment}

\paragraph{Upstream as a TurboQuant improvement, not a standalone codec.}
The $+0.37\,$dB $D_4$ shaping gain over $\Z^4$ is \emph{small}
compared to the $\sim\!1000\times$ K-MSE that TurboQuant's engineering
stack already delivers. The measured $8$--$9\%$ K-MSE improvement at
matched bits and the $1.8\times$--$5.3\times$
$\abs{\Delta\text{ppl}}$ improvement at the aggressive point are real
and architecture-independent, but they live on top of the TurboQuant
structure. The honest positioning is therefore:

\begin{quote}
\emph{v1.4 KakeyaLattice is the drop-in replacement for TurboQuant's
int$\cdot$per-coord scalar quantiser inside its Hadamard + qmax + norm
wrapper. Every TurboQuant deployment becomes ${\sim}8\%$ more faithful
on $K$ at zero engineering cost, plus the aggressive-compression
recovery.}
\end{quote}

\paragraph{Proposed upstream PR scope.}
\begin{itemize}[leftmargin=*]
  \item A Triton kernel for Conway--Sloane $D_4$ closest-point
    ($\sim\!10$ lines of element-wise ops).
  \item Slot-format alignment: $32 \times 32$ bits lattice $+\,2\cdot$fp16
    scalars per vector, bit-compatible with TQ $k8v4$.
  \item Register as a new \texttt{--kv-cache-dtype} preset, e.g.\
    \texttt{turboquant\_d4v4}, beside the existing
    \texttt{turboquant\_k8v4}.
\end{itemize}

\paragraph{Compression vs.\ per-bit fidelity.}
At every matched bit point, v1.4 pays $\sim\!3$--$11\%$ more bits than
TQ for $8$--$36\%$ better K-MSE. The \emph{Pareto curve} is strictly
better (Table~\ref{tab:pareto-qwen3}); the \emph{raw compression
ratio at a fixed operating point} is slightly worse. In deployments
where quality is the binding constraint (GLM-4-9B-Chat-class models,
long contexts with critical prompt fidelity), v1.4 is a clean win;
in raw-ratio deployments at relaxed quality, TQ's bit-level is
sufficient.

\paragraph{End-to-end vLLM throughput (not measured here).}
The snapshot harness captures $K/V$ and replaces in-place on an alt
forward, so decode-pass speed is $\sim\!$bf16 (same FlashAttention
kernel). Only the codec itself costs time. Real deployment numbers
would require a paged-cache backend using v1.4 at the kernel level;
we defer this to the upstream-PR work package.

\section{Related work}
\label{sec:related}

\paragraph{TurboQuant and the vLLM KV-cache line.}
The TurboQuant~\cite{turboquant-vllm} codec family (Hadamard rotation
+ Lloyd-Max + optional norm correction, with a $k8v4$ asymmetric
preset) was merged into vLLM on 2026-04-15 via
PR~\#38479~\cite{vllm-pr-38479}. Subsequent cleanup PRs include
\#40194~\cite{vllm-pr-40194} (Hadamard sign attribution
simplification), \#39953~\cite{vllm-pr-39953} (ROCm routing), and
\#40060~\cite{vllm-pr-40060} (CUDA backend selection). The long-context
NIAH evaluation PR~\#40122~\cite{vllm-pr-40122} confirms
$100\%$ recall at $\leq 32$k context for three of four presets.
Earlier members of the line include KIVI~\cite{kivi}
(per-channel asymmetric quantisation) and the Hadamard-based
\emph{Cache Me If You Must}~\cite{shutova-cache}.

\paragraph{Lattice-based source coding.}
Zamir and Feder's nested lattice codes~\cite{zamir-feder} provide the
capacity-achieving MSE-optimal source codec for Gaussian channels.
The $D_4$ root lattice is optimal for sphere packing in $\R^4$ and
for the second moment of its Voronoi cell
\cite{conway-sloane-book,eyuboglu-forney}. Conway and
Sloane~\cite{conway-sloane-d4} give the O($D$) closest-point
algorithm we use. Leech~$\Lambda_{24}$ would contribute $+1.53\,$dB
(vs.\ $D_4$'s $+0.37\,$dB) but requires $1000+$ lines of decoder
with severe GPU warp divergence; we judged it not worth the cost at
current noise floors.

\paragraph{Prior Kakeya-inspired KV compression.}
\cite{kakeya-v13-paper} is our prior work, introducing a PCA + WHT +
Lloyd-Max pipeline framed against the Kakeya--Brascamp--Lieb--Tropp
chain. That paper's rank cap $r=D/2$ relied on a conjectural
(Kakeya-maximal-function) lower bound; this paper's $D_4$-lattice
construction uses the same chain but relies only on the
unconditional lattice-density results of Conway--Sloane and
Zamir--Feder, making Proposition~\ref{prop:d4-kakeya} unconditional.

\paragraph{Kakeya geometry and harmonic analysis.}
The Kakeya problem~\cite{kakeya-problem} and its harmonic-analytic
role~\cite{fefferman-ball,bourgain-kakeya,wolff-kakeya,katz-laba-tao,dvir-finite-field,bennett-carbery-tao,guth-endpoint,bennett-carbery-christ-tao,carbery-valdimarsson};
the $\R^3$ resolution by Wang and Zahl~\cite{wang-zahl}.
Halko--Martinsson--Tropp~\cite{halko-martinsson-tropp} close the
chain from multilinear Kakeya through Brascamp--Lieb and matrix
Chernoff~\cite{tropp-matrix-chernoff} to numerical linear algebra.

\paragraph{KV eviction and low-rank methods.}
H$_2$O~\cite{h2o} and Scissorhands~\cite{scissorhands} are token
eviction approaches; StreamingLLM~\cite{streaming-llm} preserves
attention-sink tokens exactly; DMC~\cite{dmc} and LESS~\cite{less}
learn low-rank projections during training. All are orthogonal to
the codec line; v1.4 KakeyaLattice is a post-hoc training-free codec.

\section{Conclusion}

We presented \textsc{KakeyaLattice}~v1.4, a Zamir--Feder $D_4$
nested-lattice codec wrapped in TurboQuant's Hadamard + per-vector
qmax + unit-norm engineering stack. It is the first Kakeya-inspired
discrete-codebook construction to beat TurboQuant head-to-head on
real vLLM on four open-source model families: $12/12$ K-MSE wins,
$9/12$ $\abs{\Delta\text{ppl}}$ wins, with $128$k KV storage reduced
by $2.5\times$--$2.8\times$ at balanced quality and $1.66\times$--$1.73\times$
at near-lossless quality. The $D_4$ shaping gain of $+0.37\,$dB
over $\Z^4$ was predicted from lattice geometry and measured to
within $\sim\!1\%$ of prediction across all four models.

Three independent Kakeya-style constructions --- Guth--Katz polynomial
partitioning, raw Zamir--Feder, and empirical non-Gaussian shaping
--- failed against TurboQuant on their own. Decomposing Bridge~B's
$1414\times$ loss revealed that TurboQuant's four engineering levers
contribute $\sim$ $1000\times$ K-MSE while the $D_4$ structure
contributes only $0.92\times$: the engineering stack dominates the
codebook choice. This is the paper's central negative result and it
motivates the specific positive result of v1.4 (the engineering
stack plus a $D_4$ codebook).

The honest deployment scope is $\approx 8\%$ improvement to the
upstream TurboQuant kernel rather than a standalone Kakeya-v1.3
product. That is a real and architecture-independent gain that
applies at zero engineering cost. It closes the Kakeya-research line
in this project as a deployment search.

\paragraph{Reproducibility.}
All code, unit tests, per-model JSON measurements, snapshot hooks
for four attention architectures, and the multi-model head-to-head
harness are released under Apache-2.0. Repository:
\begin{center}
\url{https://github.com/FluffyAIcode/LLM-KV--Cache-compress}
\end{center}
The v1.4 release lives at tag \texttt{v1.4}; the canonical codec
class is \texttt{kakeyaturbo\_py.V14KakeyaZamirLatticeGPU}; the
multi-model report is at
\texttt{reports/v1\_3\_ppl/snapshot\_mode\_qwen3/multimodel/FINDINGS\_MULTIMODEL.md};
and the $128$k storage report is at
\texttt{reports/v1\_4\_release/kv\_128k\_report/V14\_KV\_128K\_REPORT.md}.

\subsection*{A note on terminology}
\label{sec:mse-discussion}

Throughout this paper we report reconstruction MSE (and its
relative normalisation) rather than inner-product distortion. For
any unit-norm query $q$,
$\abs{q^\top(k-\hat k)} \leq \norm{k-\hat k}$, so an MSE bound is a
conservative upper bound on the attention-logit perturbation that
actually matters for downstream quality. The measured
$\abs{\Delta\text{ppl}}$ and top-1 agreement corroborate that the
MSE proxy is tight in practice.

\begin{thebibliography}{99}

\bibitem{turboquant-vllm}
vibhavagarwal5 \emph{et al.}
\newblock TurboQuant: 2-bit KV cache compression with 4x capacity.
\newblock \emph{vLLM merged attention backend}, 2026.
\newblock (Merged into vllm-project/vllm main on 2026-04-15.)

\bibitem{vllm-pr-38479}
vibhavagarwal5.
\newblock [Attention Backend] TurboQuant: 2-bit KV cache compression with 4x capacity.
\newblock \emph{vllm-project/vllm Pull Request \#38479}, 2026. \url{https://github.com/vllm-project/vllm/pull/38479}.

\bibitem{vllm-pr-40194}
Alistarh, D.
\newblock [Attention] TurboQuant: remove redundant random signs, add prior art attribution.
\newblock \emph{vllm-project/vllm Pull Request \#40194}, 2026.

\bibitem{vllm-pr-39953}
aditi-amd.
\newblock [ROCm] Fix TurboQuant on ROCm: backend routing, flash-attn compat.
\newblock \emph{vllm-project/vllm Pull Request \#39953}, 2026.

\bibitem{vllm-pr-40060}
mgoin.
\newblock Fix TURBOQUANT backend selection in \texttt{cuda.py}.
\newblock \emph{vllm-project/vllm Pull Request \#40060}, 2026.

\bibitem{vllm-pr-40122}
huangzhilin-hzl.
\newblock [TurboQuant] Add NIAH long-context eval for TurboQuant KV cache.
\newblock \emph{vllm-project/vllm Pull Request \#40122}, 2026.

\bibitem{pr17-snapshot}
FluffyAIcode.
\newblock Snapshot-mode KV compression on vLLM: \texttt{reports/v1\_3\_ppl/snapshot\_mode/FINDINGS.md}.
\newblock \emph{LLM-KV--Cache-compress PR \#17}, 2026.

\bibitem{kakeya-v13-paper}
Li, A.
\newblock {Randomized Kakeya Skeletons for LLM KV Cache Compression: Algorithm and Rate--Distortion Boundary.}
\newblock \emph{LLM-KV--Cache-compress v1.3 paper}, April 2026.
\newblock (Superseded by this paper at tag \texttt{v1.4}.)

\bibitem{kivi}
Liu, Z., Yuan, J., Jin, H., Zhong, S., Xu, Z., Braverman, V., Chen, B., Hu, X.
\newblock {KIVI: A tuning-free asymmetric 2bit quantization for KV cache.}
\newblock \emph{ICML}, 2024.

\bibitem{shutova-cache}
Shutova, E. \emph{et al.}
\newblock {Cache Me If You Must: Adaptive KV Cache Quantization for LLMs.}
\newblock \emph{arXiv preprint arXiv:2501.19392}, 2025.

\bibitem{zamir-feder}
Zamir, R., Feder, M.
\newblock {On lattice quantization noise.}
\newblock \emph{IEEE Transactions on Information Theory}, 42(4):1152--1159, 1996.

\bibitem{conway-sloane-d4}
Conway, J.~H., Sloane, N.~J.~A.
\newblock {Fast quantizing and decoding algorithms for lattice quantizers and codes.}
\newblock \emph{IEEE Transactions on Information Theory}, 28(2):227--232, 1982.

\bibitem{conway-sloane-book}
Conway, J.~H., Sloane, N.~J.~A.
\newblock \emph{Sphere Packings, Lattices and Groups}, 3rd ed.
\newblock Springer, 1999.

\bibitem{eyuboglu-forney}
Eyuboglu, M.~V., Forney, G.~D.
\newblock {Lattice and trellis quantization with lattice- and trellis-bounded codebooks---high-rate theory for memoryless sources.}
\newblock \emph{IEEE Transactions on Information Theory}, 39(1):46--59, 1993.

\bibitem{h2o}
Zhang, Z., Sheng, Y., Zhou, T., Chen, T., Zheng, L., Cai, R., Song, Z., Tian, Y., R\'e, C., Barrett, C., Wang, Z., Chen, B.
\newblock {H$_2$O: Heavy-hitter oracle for efficient generative inference of large language models.}
\newblock \emph{NeurIPS}, 2023.

\bibitem{scissorhands}
Liu, Z., Desai, A., Liao, F., Wang, W., Xie, V., Xu, Z., Kyrillidis, A., Shrivastava, A.
\newblock {Scissorhands: Exploiting the persistence of importance hypothesis for LLM KV cache compression at test time.}
\newblock \emph{NeurIPS}, 2023.

\bibitem{streaming-llm}
Xiao, G., Tian, Y., Chen, B., Han, S., Lewis, M.
\newblock {Efficient streaming language models with attention sinks.}
\newblock \emph{ICLR}, 2024.

\bibitem{dmc}
Nawrot, P., Lancucki, A., Chochowski, M., Tarjan, D., Ponti, E.~M.
\newblock {Dynamic memory compression: Retrofitting LLMs for accelerated inference.}
\newblock \emph{arXiv preprint arXiv:2403.09636}, 2024.

\bibitem{less}
Dong, H., Yang, X., Zhang, Z., Wang, Z., Chi, Y., Chen, B.
\newblock {Get more with LESS: Synthesizing recurrence with KV cache compression for efficient LLM inference.}
\newblock \emph{ICML}, 2024.

\bibitem{deepseekv3}
DeepSeek-AI.
\newblock {DeepSeek-V3 Technical Report.}
\newblock \emph{arXiv preprint arXiv:2412.19437}, 2024.

\bibitem{hadamard-bias}
Chee, J., Cai, Y., Kuleshov, V., De~Sa, C.
\newblock {QuIP\#: Even better LLM quantization with hadamard incoherence and lattice codebooks.}
\newblock \emph{ICML}, 2024.

\bibitem{kakeya-problem}
Kakeya, S.
\newblock {Some problems on maxima and minima regarding ovals.}
\newblock \emph{Tohoku Science Reports}, 6:71--88, 1917.

\bibitem{fefferman-ball}
Fefferman, C.
\newblock {The multiplier problem for the ball.}
\newblock \emph{Annals of Mathematics}, 94(2):330--336, 1971.

\bibitem{bourgain-kakeya}
Bourgain, J.
\newblock {Besicovitch type maximal operators and applications to Fourier analysis.}
\newblock \emph{Geometric and Functional Analysis}, 1(2):147--187, 1991.

\bibitem{wolff-kakeya}
Wolff, T.
\newblock {An improved bound for Kakeya type maximal functions.}
\newblock \emph{Revista Matem\'atica Iberoamericana}, 11(3):651--674, 1995.

\bibitem{katz-laba-tao}
Katz, N.~H., {\L}aba, I., Tao, T.
\newblock {An improved bound on the Minkowski dimension of Besicovitch sets in $\R^3$.}
\newblock \emph{Annals of Mathematics}, 152(2):383--446, 2000.

\bibitem{dvir-finite-field}
Dvir, Z.
\newblock {On the size of Kakeya sets in finite fields.}
\newblock \emph{Journal of the American Mathematical Society}, 22(4):1093--1097, 2009.

\bibitem{bennett-carbery-tao}
Bennett, J., Carbery, A., Tao, T.
\newblock {On the multilinear restriction and Kakeya conjectures.}
\newblock \emph{Acta Mathematica}, 196(2):261--302, 2006.

\bibitem{guth-endpoint}
Guth, L.
\newblock {The endpoint case of the Bennett--Carbery--Tao multilinear Kakeya conjecture.}
\newblock \emph{Acta Mathematica}, 205(2):263--286, 2010.

\bibitem{guth-katz}
Guth, L., Katz, N.~H.
\newblock {On the Erd{\H o}s distinct distances problem in the plane.}
\newblock \emph{Annals of Mathematics}, 181(1):155--190, 2015.

\bibitem{bennett-carbery-christ-tao}
Bennett, J., Carbery, A., Christ, M., Tao, T.
\newblock {The Brascamp--Lieb inequalities: finiteness, structure and extremals.}
\newblock \emph{Geometric and Functional Analysis}, 17(5):1343--1415, 2008.

\bibitem{carbery-valdimarsson}
Carbery, A., Valdimarsson, S.~I.
\newblock {The endpoint multilinear Kakeya theorem via the Borsuk--Ulam theorem.}
\newblock \emph{Journal of Functional Analysis}, 264(7):1643--1663, 2013.

\bibitem{tropp-matrix-chernoff}
Tropp, J.~A.
\newblock {User-friendly tail bounds for sums of random matrices.}
\newblock \emph{Foundations of Computational Mathematics}, 12(4):389--434, 2012.

\bibitem{halko-martinsson-tropp}
Halko, N., Martinsson, P.-G., Tropp, J.~A.
\newblock {Finding structure with randomness: Probabilistic algorithms for constructing approximate matrix decompositions.}
\newblock \emph{SIAM Review}, 53(2):217--288, 2011.

\bibitem{wang-zahl}
Wang, H., Zahl, J.
\newblock {Volume estimates for unions of convex sets, and the Kakeya set conjecture in three dimensions.}
\newblock \emph{arXiv preprint arXiv:2502.17655}, 2025.

\end{thebibliography}

\appendix

\section{Reproducibility manifest}
\label{app:repro}

The full multi-model benchmark is reproducible via:
\begin{verbatim}
cd /workspace/LLM-KV--Cache-compress
export VLLM_ENABLE_V1_MULTIPROCESSING=0 KAKEYA_SNAPSHOT_QWEN3=1
python benchmarks/multimodel_v14_vs_tq.py \
    --model-path Qwen/Qwen3-4B --model-name qwen3_4b \
    --ctx-len 2048 --n-eval 64 --n-passages 4 \
    --out-dir reports/v1_3_ppl/snapshot_mode_qwen3/multimodel
\end{verbatim}
Replace \texttt{--model-path} / \texttt{--model-name} for the other
three models; add \texttt{--trust-remote-code} for GLM-4-9B-Chat.
The $128$k storage report uses
\texttt{benchmarks/multimodel\_v14\_kv\_128k\_report.py} with the
same capture/replace protocol and both $K$ and $V$ compressed.

\section{Canonical naming}
\label{app:naming}

\begin{itemize}[leftmargin=*]
  \item \textbf{Spoken / written name}: ``v1.4 kakeya zamir lattice GPU''
  \item \textbf{Short form with parameter}: e.g.\ ``v1.4 Q=152''
  \item \textbf{Class name}: \texttt{V14KakeyaZamirLatticeGPU}
  \item \textbf{Module}: \texttt{kakeyaturbo\_py.v1\_4\_kakeya\_zamir\_lattice\_gpu}
\end{itemize}
The research-lineage name ``Bridge B2'' is reserved for the research
prototype module \texttt{bridge\_b2\_d4\_tq\_style.py} and the
session's git-blame history; new writing must use the v1.4 name.

\section{Canonical operating points}
\label{app:ops}

\begin{tabular}{rrrl}
  \toprule
  $q_\mathrm{range}$ & bits/tok/head & CR (bf16) & Role \\
  \midrule
  $152$ & $1088$ & $1.88\times$ & Near-lossless, head-to-head with TQ $k8v4$ \\
  $\phantom{0}76$ & $\phantom{0}960$ & $2.13\times$ & \\
  $\phantom{0}38$ & $\phantom{0}832$ & $2.46\times$ & \textbf{Deployment recommended balanced} \\
  $\phantom{0}19$ & $\phantom{0}704$ & $2.91\times$ & \\
  $\phantom{0}10$ & $\phantom{0}576$ & $3.56\times$ & Aggressive compression \\
  $\phantom{00}5$ & $\phantom{0}448$ & $4.57\times$ & \\
  $\phantom{00}2$ & $\phantom{0}320$ & $6.40\times$ & Extreme compression \\
  \bottomrule
\end{tabular}

\section{Failed alternatives catalogue}
\label{app:failed}

For audit-trail completeness, beyond Bridges A/B/C
(Section~\ref{sec:bridges}), the following standalone constructions
were also tested and rejected during the project's lifetime:

\begin{enumerate}[leftmargin=*]
  \item \textbf{Besicovitch skeleton replacing PCA on $K$.}
    Loses $18$\,pp $\Delta$ppl at matched MSE. The Haar codebook's
    rotation-invariance defeats attention-weighted rotation on $K$.
  \item \textbf{K-residual Besicovitch.} $3\times$ worse MSE than
    Lloyd-Max at matched bits; WHT + Lloyd-Max is a joint design.
  \item \textbf{Perron-tree / attention-weighted direction density.}
    Oracle MSE gain $0.1\%$ on real data; Haar codebook is already
    near-optimal on isotropic $V$.
  \item \textbf{Non-Gaussian shaping (Bridge~C, this paper
    Section~\ref{sec:bridge-c}).} Finite-sample noise in empirical
    Lloyd-Max centroids outweighs the $\sim\!2\times$ non-Gaussian
    headroom.
  \item \textbf{$8$-layer boundary expansion.} Strictly worse than
    $6$-layer at identical recipe; extra protected layers leave
    more compressed layers at worse Q.
  \item \textbf{Leech lattice $\Lambda_{24}$.} $+1.53\,$dB shaping
    gain but $50$--$100\times$ slower encode on GPU (warp divergence
    in the Conway decoder). Not worth pursuing at current noise
    floors.
\end{enumerate}

\end{document}
